\section{What is being counted, and what is proved internally}

The number is \textbf{five}, where position, orientation, and positive scale
are disregarded. There are of course infinitely many placed or scaled
tetrahedra. Reflections are allowed in the equivalence relation.

\begin{definition}[Geometric input]
A Platonic polyhedron is a bounded, full-dimensional convex polyhedron
whose maximal two-dimensional faces are congruent convex regular polygons,
and for which the same number of faces meets at every vertex. If each
face has $p$ sides and each vertex has valence $q$, its type is $(p,q)$.
Faces are maximal facets: an artificial subdivision does not create new
faces for this definition.
\end{definition}

We retain this usual local definition. In particular we do not assume
flag transitivity, or assume that the type already determines the shape.
Those shortcuts would hide part of the classification problem.

\subsection{The finite coordinate domain}

First-order arithmetic has natural numbers as its objects. A sentence
about all unspecified set-theoretic real polyhedra is not literally a
sentence of $\ipra$. Our internal objects therefore have the following
explicit representation.

A scalar is a finite expression obtained from rational numbers by
addition, subtraction, multiplication, division by a nonzero scalar,
and the nonnegative square root of a nonnegative scalar. Equivalently,
place a finite collection of such expressions in a common tower
\[
 \Q=K_0\subseteq K_1\subseteq\cdots\subseteq K_n,
 \qquad K_j=K_{j-1}(s_j),\quad s_j^2=r_j\geq0.
\]
The extensions need not have degree two: $r_j$ may be zero or a square
already in $K_{j-1}$. This redundancy must be respected by equality.
All points of the input use one common finite tower. New auxiliary square
roots are appended to it. This is the field of real constructible scalars,
presented by finite codes; it is not a representation of arbitrary
computable real numbers.

A polyhedron code lists distinct vertices, its facet cycles, and their
incidences. Face cycles have distinct vertices, and duplicate facet
vertex sets are forbidden. The validator checks the following finite conditions:
\begin{enumerate}
\item Four listed vertices are affinely independent. Each listed vertex
is an extreme point, and every facet is a convex planar cycle.
\item For each facet an oriented cross-product normal gives a supporting
plane. Its equality set among the vertices is exactly the facet's vertex
set; all other vertices lie strictly on the inner side.
\item All facets occur. This is checked by examining every noncollinear
vertex triple: whenever its plane is supporting, its full equality set
must be one of the listed facets.
\item Edges are the consecutive vertex pairs on facet boundaries.
Every edge belongs to two facets; each vertex link is one cyclic list;
the face adjacency graph is connected. These are also consequences of
the convex hull description, but checking them makes the data interface
particularly transparent.
\item The facets have common size $p\geq3$, are regular and congruent,
and the vertex links have common length $q\geq3$.
\end{enumerate}
Extreme-point status can be checked without linear programming: the
normals of the incident supporting planes must span three dimensions.
For a supporting-plane intersection this is equivalent to a zero-dimensional
face. Completeness of the facet list follows from the elementary fact
that every facet of a full-dimensional finite hull contains three
affinely independent vertices. Conversely the equality set of every
supporting noncollinear triple is exactly the vertex set of a facet,
so the test adds no extraneous facets.

Regularity has a convenient exact test. Center the ordered vertices
$x_0,\ldots,x_{p-1}$ of a face at their average. Require its Gram matrix
to be invariant under $(i,j)\mapsto(i+1,j+1)$ modulo $p$. Cyclic shift
then extends to an isometry of the face plane. Since it preserves the
convex cyclic order, it is a rotation carrying each vertex to its
successor. This is exactly a convex regular polygon. Equality of the
corresponding face Gram matrices tests congruence. Thus no trigonometric
function is needed even to state regularity.

\begin{theorem}[Internal finite classification]\label{thm:main}
There are fixed primitive recursive functions $\Valid$, $\Class$, $\Witness$,
and $\Similar$, and five fixed model codes $\Model(0),\ldots,\Model(4)$,
for which $\ipra$ proves
\begin{align*}
 \Valid(c)=1\ \Longrightarrow\ &\Class(c)<5\ \wedge\\[-4pt]
 &\Similar(c,\Model(\Class(c)),\Witness(c))=1.
\end{align*}
It also proves that the five model codes are valid and pairwise
non-similar. The similarity certificate contains a vertex bijection
and an affine map $x\mapsto Ax+b$ satisfying
$A^{\mathsf T}A=\lambda I$ for a positive scalar $\lambda$.
\end{theorem}

This is a theorem with a constructive witness, rather than an appeal to
a quotient whose number of elements is assumed to exist. The proof of
the theorem is the mathematical argument below, together with its
explicit primitive recursive implementation described in
Section~\ref{sec:internal}. A machine-checked derivation in the embedded
syntax remains a formalization task.

\begin{remark}[Relation to the ordinary theorem]
All five representatives have coordinates in $\Q(\sqrt5)$. The same
geometric proof, interpreted externally over ordinary $\R$, gives the
usual unrestricted classical classification. Thus the finite theorem
captures all five shapes. It does not assert that arbitrary real-number
names can be compared or normalized by a primitive recursive program.
\end{remark}

\subsection{Exact order and why the scalar algorithms are primitive recursive}

The only potentially delicate computational ingredient is exact sign.
Here are explicit algorithms, including redundant tower stages.
At one stage write $x=a+bs$, where $s^2=r\geq0$ and all of $a,b,r$
belong to the preceding stage. If $r=0$, return the sign of $a$.
Otherwise handle $a=0$ or $b=0$ directly. If $a,b$ have the same sign,
that is the sign of $x$. If their signs are opposite, then
\begin{equation}\label{eq:towersign}
 \sgn(a+b\sqrt r)=\sgn(a)\,\sgn(a^2-b^2r).
\end{equation}
Indeed, when $a>0>b$, comparison with zero is precisely the comparison
of the two positive quantities $a$ and $(-b)\sqrt r$; squaring preserves
that comparison. The other case follows by negation. Every comparison
on the right is in a strictly shorter tower.

Addition and multiplication are
\[
 (a,b)+(c,d)=(a+c,b+d),\qquad
 (a,b)(c,d)=(ac+bd r,ad+bc).
\]
For inversion put $N=a^2-b^2r$. If $N\ne0$, use
\[
 (a+bs)^{-1}=(a-bs)/N.
\]
If $N=0$ and $r>0$, then $a$ and $bs$ have equal absolute value.
If $a=0$ or the signs of $a,b$ are opposite, the element is zero;
otherwise $a+bs=2a\ne0$, and its inverse is $1/(2a)$ in the preceding
stage. For $r=0$ invert $a$ directly. Define the inverse of zero to be
zero as a totalization, and reject a division by zero when validating
an expression. Equality is sign zero, \emph{not} coefficientwise equality
at a possibly redundant stage. The following lemma proves the required
ordered-field laws; quotient notation is only shorthand for the
decidable equality relation on codes.

\begin{lemma}[One ordered quadratic stage]\label{lem:orderedstage}
Let $K$ be a decidably ordered field and $r\geq0$ in $K$. On formal
pairs $a+bs$ use the displayed addition, multiplication, and sign
algorithms, with $s^2=r$. Define $x\sim y$ when $\sgn(x-y)=0$.
Then $\sim$ is a ring congruence, sign is well-defined modulo $\sim$,
and the represented structure is an ordered field with $s\geq0$
and $s^2=r$.
\end{lemma}
\begin{proof}
If $r=0$, operations and sign reduce to evaluation $a+bs\mapsto a$.
Assume $r>0$. First, whenever $t>0$ in $K$ satisfies $t^2=r$, the
prescribed sign of every $c+ds$ equals $\sgn_K(c+dt)$. In the
opposite-sign case this is precisely comparison by squaring the
two nonnegative quantities; all other cases are immediate.
Evaluation at such a $t$ respects addition and multiplication.

This observation handles redundant representations without a test
for whether $r$ has a square root in $K$. If $\sgn(a+bs)=0$ but
$(a,b)\ne(0,0)$, the sign cases give $b\ne0$ and the explicit root
$t=-a/b>0$, $t^2=r$. Whenever an equality premise has a nontrivial
zero difference, this $t$ reduces sign invariance, transitivity,
and compatibility with operations to the corresponding facts in $K$.
If every zero difference in the premises is the literal coefficient-zero
pair, the facts hold coefficientwise. Hence $\sim$ is an equivalence
and ring congruence, and sign is invariant under $\sim$.
The ring identities themselves follow by expanding coefficients
using $s^2=r$. Directly from the sign rule,
$\sgn(-x)=-\sgn(x)$, and multiplication by a positive scalar from
$K$ preserves sign.

We explicitly prove the two order-closure properties. Define predicates
solely in the base field:
\[
 L(u)\ \Longleftrightarrow\ u<0\ \text{or}\ (u\geq0\ \text{and}\ u^2<r),
 \qquad
 U(u)\ \Longleftrightarrow\ u>0\ \text{and}\ u^2>r.
\]
$L$ is downward closed, $U$ is upward closed, and $L(l),U(u)$ imply
$l<u$, by monotonicity of squaring on nonnegative elements.
Every positive pair has exactly one of these three forms: a positive
scalar; $b(s-l)$ with $b>0$ and $L(l)$; or $e(u-s)$ with $e>0$ and
$U(u)$. This rewrites the sign rule according to the sign of the
$s$ coefficient.

Two lower forms add to a lower form whose threshold is the positive
weighted average of their thresholds. This is no greater than their
maximum, so $L$ holds. Two upper forms use the corresponding minimum
and $U$. Adding a positive scalar lowers a lower threshold or raises
an upper threshold. For the remaining case add
$b(s-l)+e(u-s)$, where $l<u$. If $b=e$ the sum is $b(u-l)>0$.
If $b>e$, it is $(b-e)(s-l')$, with
\[
 l'=l-\frac{e(u-l)}{b-e}<l,
\]
so $L(l')$ holds. If $b<e$, it is $(e-b)(u'-s)$, with
\[
 u'=u+\frac{b(u-l)}{e-b}>u,
\]
so $U(u')$ holds. Thus the sum of two positive pairs is positive,
without presupposing an order on the extension.

For multiplication put $N(a,b)=a^2-b^2r$. Expansion gives
$N(xy)=N(x)N(y)$. If a positive $x=a+bs$ has $N(x)=0$, its
coefficients $a,b$ are both positive. The explicit root $t=a/b$
and the evaluation observation reduce positivity of its product
with any positive pair to positivity in $K$. The same applies to $y$.
Otherwise, positivity and $N(x)>0$ imply $a>0$, $a^2>b^2r$;
positivity and $N(x)<0$ imply $b>0$, $b^2r>a^2$.
Write $y=c+ds$ and $xy=A+Bs$, where
$A=ac+bdr$, $B=ad+bc$. The four possibilities are
\[
\begin{array}{ccl}
 N(x)>0,\ N(y)>0 &: & ac>|bd|r,\quad A>0,\quad N(xy)>0;\\
 N(x)<0,\ N(y)<0 &: & bdr>|ac|,\quad A>0,\quad N(xy)>0;\\
 N(x)>0,\ N(y)<0 &: & ad>|bc|,\quad B>0,\quad N(xy)<0;\\
 N(x)<0,\ N(y)>0 &: & bc>|ad|,\quad B>0,\quad N(xy)<0.
\end{array}
\]
Each first inequality follows by squaring its positive left side
and multiplying the preceding nonnegative inequalities, cancelling
$r>0$ in the mixed cases. For example, in the third case
\[
 (a^2-b^2r)d^2r+b^2r(d^2r-c^2)>0
\]
is exactly $r((ad)^2-(bc)^2)>0$. The first summand is strictly
positive, since $d>0$. The sign algorithm therefore declares the
product positive in all four cases.

Sign trichotomy, negation, congruence, and the two closure results
give a translation-invariant linear order compatible with
multiplication. For inverses, when $N(x)\ne0$ the conjugate formula
multiplies to one coefficientwise. When $N(x)=0$ but $x$ is nonzero
modulo $\sim$, $a,b$ have the same nonzero sign; $t=a/b>0$ satisfies
$t^2=r$, evaluation sends $x$ to $2a$, and the scalar $1/(2a)$
multiplies with $x$ to one modulo $\sim$. This proves the field
assertion. The element $(0,1)$ is positive and squares to $(r,0)$.

This proof supplies a fixed finite family of base-field templates:
the equivalence and congruence laws, commutative-ring identities,
sign trichotomy and negation, closure of positivity under addition
and multiplication, and the nonzero-inverse identity. Every division
has an explicit nonzero premise; every split is a decidable scalar
comparison. These are the templates used in the bounded law-table
induction below.
\end{proof}

For clarity, a size bound accompanies this recursive description.
Store a level-$n$ element as a dense binary tree of rational coefficient
pairs, and let $H\geq1$ bound the binary lengths of all input and tower
radicand numerators and denominators. Addition, multiplication, sign,
and the displayed inversion algorithm make at most sixteen calls to
strictly lower-level routines at any node. Scalar multiplication here
is a lower-level multiplication, not an additional same-level call.
There are therefore fewer than $16^{n+1}$ nodes in an operation's
expanded call tree. A deliberately larger bound, allowing validation
and bookkeeping, is
\[
 L(n)=2^{16(n+1)}.
\]
Unnormalized rational operations change a coefficient length bound $B$
to at most $2B+1$. Consequently
\begin{equation}\label{eq:bitbound}
 B(n,H)=(H+2)\,2^{L(n)}
\end{equation}
bounds the intermediate coefficient lengths. Standard binary coding of
the finite coefficient arrays and an explicit stack now implement the
algorithms by bounded loops. Long integer arithmetic also has bounded
loops. Both bounds are primitive recursive. Finite sequences of these
operations use primitive recursion on their length, updating the bound
at each step. Appending a root increases the tower depth by one.

These are bounds for computations of the basic operations. A sequence
of $K$ operations is bounded by replacing $L(n)$ with $K L(n)$, or by
iterating the one-operation bound. The ordered-field laws for
represented equality require the bounded law-table construction in
Section~\ref{sec:internal}: fix the final height and coefficient budget,
propagate larger budgets to shorter prefixes, and induct on equations
asserting correctness of the finite tables. Thus no induction on an
unbounded universally quantified field-law assertion is required.
Both the size bounds and the correctness argument matter; mere
termination of an unspecified symbolic procedure would not suffice.

\subsection{The arithmetic theory being used}

We use intuitionistic first-order logic with equality, symbols for
primitive recursive functions, their defining equations, the successor
axioms, and induction for quantifier-free formulas. This is the
presentation in ProveIt's \code{ConciseIPRA/PR.lean}~\cite{proveit}.
We do not add induction for arbitrary formulas.

A bounded assertion over a coded list is represented by a primitive
recursive Boolean fold. An induction invariant is an equation saying
that this fold equals one, together with equations for the accumulated
counts and witnesses. Thus bounded universal checks do not require
stronger induction. All case splits below are on natural-number or exact
scalar comparisons whose decisions were just constructed. We use neither
excluded middle for arbitrary propositions, Markov's principle, choice,
nor a universal evaluator for arbitrary primitive recursive programs.
